% ===== MAIN.TEX =====
% Entry point cho báo cáo LaTeX
% Compile: xelatex main.tex && biber main && xelatex main.tex && xelatex main.tex

% ====== PREAMBLE.TEX ======
% XeLaTeX preamble cho báo cáo HUST BTL/Đồ Án
% Compile: xelatex -interaction=nonstopmode main.tex && biber main && xelatex main.tex

\documentclass[12pt,a4paper,oneside]{book}

% ===== FONT & LANGUAGE =====
\usepackage{fontspec}
\usepackage{polyglossia}
\setdefaultlanguage{vietnamese}
\setotherlanguage{english}
\usepackage{lmodern}

% Thiết lập font chính (Times New Roman 12pt)
% Nếu không có TNR, thay vào Liberation Serif hoặc TeX Gyre Termes
\setmainfont[
  Extension = .ttf,
  BoldFont = *b,
  ItalicFont = *i,
  BoldItalicFont = *bi
]{times}
% Fallback (uncomment nếu không có file font local):
% \setmainfont{Times New Roman}[SizeFeatures={Size=12}]

\setsansfont{DejaVu Sans}
\setmonofont{DejaVu Sans Mono}

% ===== PAGE LAYOUT =====
\usepackage[
  a4paper,
  margin=2.5cm,
  top=2.5cm,
  bottom=2.0cm,
  left=2.5cm,
  right=2.5cm,
  headheight=14pt,
  includehead,
  includefoot
]{geometry}

\usepackage{setspace}
\onehalfspacing % Dãn dòng 1.5

% ===== GRAPHICS & TABLES =====
\usepackage{graphicx}
\graphicspath{{figures/}}
\usepackage{float}
\usepackage{caption}
\usepackage{subcaption}
\usepackage{booktabs}
\usepackage{longtable}
\usepackage{array}
\usepackage{colortbl}

% ===== COLOR & LINKS =====
\usepackage{xcolor}
\usepackage[
  colorlinks=true,
  linkcolor=blue,
  urlcolor=blue,
  citecolor=red,
  pdfstartview=FitH,
  pdftitle={BTL/Đồ Án AI},
  pdfauthor={HUST Student}
]{hyperref}

% ===== CODE LISTINGS =====
\usepackage{listings}
\lstset{
  basicstyle=\ttfamily\small,
  breaklines=true,
  frame=single,
  framerule=0.5pt,
  numbers=left,
  numberstyle=\tiny\color{gray},
  keywordstyle=\color{blue}\bfseries,
  commentstyle=\color{gray}\itshape,
  stringstyle=\color{red!60!black},
  showstringspaces=false,
  tabsize=2,
  xleftmargin=15pt,
  framexleftmargin=10pt,
  captionpos=b,
  belowcaptionskip=8pt
}

\lstdefinelanguage{Python}{
  morekeywords={def,class,import,from,return,if,else,elif,for,while,try,except,finally,with,as,lambda,yield,pass,break,continue,raise},
  sensitive=true,
  morecomment=[l]{\#},
  morestring=[b]",
  morestring=[b]'
}

% ===== SECTION HEADING STYLE =====
\usepackage{titlesec}
\titleformat{\chapter}[display]
  {\normalfont\huge\bfseries}
  {Chương \thechapter}
  {20pt}
  {\Huge}
\titlespacing*{\chapter}{0pt}{-30pt}{20pt}

\titleformat{\section}[hang]
  {\normalfont\Large\bfseries}
  {\thesection}{12pt}{}
\titleformat{\subsection}[hang]
  {\normalfont\large\bfseries}
  {\thesubsection}{12pt}{}

% ===== TOC DEPTH & FORMATTING =====
\setcounter{tocdepth}{2}
\setcounter{secnumdepth}{3}

% ===== BIBLIOGRAPHY (biber + numeric style) =====
\usepackage[
  backend=biber,
  style=numeric,
  sorting=nyt,
  language=english
]{biblatex}
\addbibresource{refs.bib}

% ===== HEADERS & FOOTERS =====
\usepackage{fancyhdr}
\pagestyle{fancy}
\fancyhf{}
\fancyhead[L]{BTL/Đồ Án môn Nhập môn AI}
\fancyhead[R]{\thepage}
\fancyfoot[C]{\thepage}
\renewcommand{\headrulewidth}{0.5pt}
\renewcommand{\footrulewidth}{0pt}

% ===== METADATA =====
\title{Báo Cáo Bài Tập Lớn / Đồ Án}
\author{Sinh viên HUST}
\date{\today}


\begin{document}

% ===== TITLEPAGE.TEX =====
% Trang bìa với placeholder tên sinh viên
% Hướng dẫn: Sửa các lệnh \newcommand dưới đây để điền tên/MSSV/lớp/GV

\newcommand{\svname}{HỌ VÀ TÊN SINH VIÊN}
\newcommand{\svid}{MSSV}
\newcommand{\svclass}{LỚP}
\newcommand{\gvname}{TS. TÊN GIẢNG VIÊN}
\newcommand{\reporttitle}{Ứng dụng Chỉ Đường Metro Thượng Hải dùng Thuật toán Tìm Kiếm}
\newcommand{\reportsubtitle}{Phân Tích, Thiết Kế và Thực Hiện Bốn Thuật toán Pathfinding}

\begin{titlepage}
  \centering
  \vspace*{1.5cm}
  {\Large\textbf{ĐẠI HỌC BÁCH KHOA HÀ NỘI}}\\
  \vspace{0.3cm}
  {\large Khoa Công Nghệ Thông Tin}\\
  \vspace{2cm}

  {\LARGE\textbf{BÁO CÁO BÀI TẬP LỚN}}\\
  \vspace{0.5cm}
  {\large \textit{(Đồ Án môn học)}}\\
  \vspace{2cm}

  {\huge\textbf{\reporttitle}}\\
  \vspace{0.5cm}
  {\large \textit{\reportsubtitle}}\\
  \vspace{3cm}

  {\Large Môn: Nhập môn Trí Tuệ Nhân Tạo (Introduction to AI)}\\
  \vspace{2cm}

  \begin{tabular}{ll}
    \textbf{Sinh viên:} & \svname \\
    \textbf{MSSV:} & \svid \\
    \textbf{Lớp:} & \svclass \\
  \end{tabular}
  \vspace{2cm}

  {\Large\textbf{Giáo viên hướng dẫn:}}\\
  {\Large \gvname}\\
  \vspace{3cm}

  {\Large Hà Nội, \today}
\end{titlepage}

\cleardoublepage


% ===== FRONT MATTER =====
\frontmatter

\tableofcontents
\cleardoublepage

\listoffigures
\cleardoublepage

\listoftables
\cleardoublepage

% ===== LỜI CAM ĐOAN =====
\chapter*{Lời Cam Đoan}

Em xin cam đoan rằng báo cáo này là kết quả công việc nghiên cứu độc lập của em, dựa trên kiến thức được học tập từ các bài giảng và tài liệu tham khảo. Toàn bộ nội dung báo cáo đều là sáng tạo của em, không trùng lặp hay sao chép từ các công trình khác mà không dẫn chiếu.

Các điểm được trích dẫn từ công trình của người khác đều có ghi rõ nguồn tài liệu.

\vspace{2cm}

Hà Nội, \today

\hspace{3cm} Sinh viên

\vspace{2cm}

\hspace{3cm} \svname

\cleardoublepage

% ===== LỜI CẢM ƠN =====
\chapter*{Lời Cảm Ơn}

Tác giả xin cảm ơn \gvname{}, giáo viên hướng dẫn, vì những lời khuyên quý báu, sự hỗ trợ tận tình trong quá trình thực hiện bài tập lớn này.

Cảm ơn các giáo viên khác trong phòng Công Nghệ Thông Tin đã cung cấp tài liệu tham khảo và môi trường học tập tốt.

Cảm ơn các bạn cùng lớp đã trao đổi ý kiến, giúp đỡ em hoàn thiện công trình.

\cleardoublepage

% ===== TÓM TẮT =====
\chapter*{Tóm Tắt}

Báo cáo này trình bày một ứng dụng web chỉ đường trên mạng lưới tàu điện ngầm Thượng Hải, sử dụng bốn thuật toán tìm kiếm đường ngắn nhất cổ điển: Uniform Cost Search (UCS), Greedy Best-First Search, A* Search, và Dijkstra.

Báo cáo gồm 5 chương chính:

\begin{enumerate}
  \item \textbf{Mở đầu}: Nêu lý do chọn đề tài, mục tiêu nghiên cứu, phạm vi và phương pháp thực hiện.
  \item \textbf{Cơ sở lý thuyết}: Trình bày chi tiết bốn thuật toán pathfinding, hàm heuristic admissible, và so sánh lý thuyết các thuật toán.
  \item \textbf{Phân tích và thiết kế}: Phân tích yêu cầu, kiến trúc hệ thống 3-tầng, mô hình dữ liệu, thiết kế API và giao diện người dùng.
  \item \textbf{Thực hiện và kết quả}: Mô tả công nghệ sử dụng, cấu trúc mã nguồn, code listings các thuật toán, kết quả kiểm thử và benchmark so sánh.
  \item \textbf{Kết luận}: Tóm tắt những gì đạt được, hạn chế hiện tại, và hướng phát triển tương lai.
\end{enumerate}

Ứng dụng được xây dựng bằng Python (backend) và JavaScript (frontend), sử dụng dữ liệu thật từ OpenStreetMap (411 ga, 20 tuyến). Bộ test bao gồm 29 test case (unit + integration), đạt tỉ lệ coverage cao. Kết quả benchmark cho thấy A* là nhanh nhất (64 node explored), Greedy nhanh nhất về số node (15 node), và UCS, Dijkstra đều tối ưu nhưng Dijkstra tính toàn bộ SSSP nên explored nhiều node hơn.

Từ khóa: Pathfinding, Graph algorithms, A*, Dijkstra, Shanghai Metro, Web application.

\cleardoublepage

% ===== MAIN MATTER =====
\mainmatter

% ===== CHƯƠNG 1: MỞ ĐẦU =====

\chapter{Mở Đầu}

\section{Lý do chọn đề tài}

Bài toán tìm đường ngắn nhất trên đồ thị có trọng số là một trong những vấn đề cơ bản nhất của Khoa học Máy tính và Trí tuệ Nhân tạo, với ứng dụng rộng rãi trong đời sống thực tế. Các thuật toán pathfinding được sử dụng trong hệ thống định vị GPS, lập kế hoạch tuyến đường giao thông, trò chơi điện tử, robot navigation, và nhiều lĩnh vực khác.

Thành phố Thượng Hải sở hữu một trong những hệ thống tàu điện ngầm (Metro) lớn nhất thế giới với hơn 400 ga và 20 tuyến, tạo ra một đồ thị phức tạp lý tưởng để kiểm thử các thuật toán. Việc xây dựng một ứng dụng thực tế để tìm đường tối ưu giữa hai ga metro không chỉ giúp hiểu rõ lý thuyết mà còn minh hoạ cách áp dụng các thuật toán vào bài toán thực tiễn.

Ngoài ra, đề tài này cho phép so sánh trực tiếp bốn thuật toán tìm kiếm khác nhau --- Uniform Cost Search, Greedy Best-First Search, A*, và Dijkstra --- trên cùng một bộ dữ liệu, từ đó rút ra các kết luận về hiệu năng, tính tối ưu, và sự phù hợp của mỗi thuật toán trong các tình huống khác nhau.

\section{Mục tiêu của đề tài}

Báo cáo này có các mục tiêu chính sau:

\begin{enumerate}
  \item \textbf{Nghiên cứu và phân tích} bốn thuật toán tìm kiếm: hiểu rõ nguyên lý hoạt động, độ phức tạp thời gian và không gian, tính hoàn chỉnh (completeness) và tính tối ưu (optimality) của từng thuật toán.

  \item \textbf{So sánh hiệu năng} của các thuật toán trên dữ liệu thực: đo lường số node được khám phá, chi phí đường đi, thời gian thực thi để rút ra những phát hiện về khi nào nên dùng thuật toán nào.

  \item \textbf{Xây dựng ứng dụng web thực tế} cho phép người dùng tìm đường giữa hai ga metro theo nhiều chế độ khác nhau, bao gồm chọn từ danh sách và click trực tiếp trên bản đồ.

  \item \textbf{Phát triển kỹ năng lập trình} về kiến trúc backend-frontend, xử lý yêu cầu HTTP, trực quan hoá dữ liệu trên bản đồ, và viết test case toàn diện.

  \item \textbf{Làm việc với dữ liệu thực} từ OpenStreetMap, hiểu những thách thức của xử lý dữ liệu địa lý, và xử lý các edge case như dữ liệu lỗi, điểm ngoài phạm vi, v.v.
\end{enumerate}

\section{Phạm vi đề tài}

Báo cáo này có các giới hạn và giả định sau:

\begin{itemize}
  \item \textbf{Dữ liệu:} Sử dụng dữ liệu mạng lưới metro Thượng Hải từ OpenStreetMap gồm 411 ga và 20 tuyến. Đây là dữ liệu sẵn có từ ngày cập nhật, không được cập nhật real-time.

  \item \textbf{Đồ thị:} Đồ thị là vô hướng (undirected), có trọng số dương, biểu diễn mạng metro nơi có thể đi cả hai chiều giữa các ga liền kề.

  \item \textbf{Chi phí cạnh:} Chi phí được ước lượng dựa trên công thức (khoảng cách / 35 km/h + phụ phí dừng), không dùng dữ liệu thời gian chạy thực tế do OSM không cung cấp.

  \item \textbf{Giới hạn địa lý:} Khi tìm ga gần nhất từ một điểm bất kỳ, hệ thống chỉ chấp nhận những điểm nằm trong bán kính 50 km từ ga gần nhất, nếu vượt quá sẽ báo lỗi.

  \item \textbf{Không tính:} Báo cáo không xét các yếu tố như giờ cao điểm, đóng cửa tuyến, độ trễ thực tế, hoặc thời gian đợi giữa các tuyến (transfer time). Đây là những yếu tố để mở rộng trong tương lai.

  \item \textbf{Một cặp (start, goal)} trong mỗi lần chạy. Không xét bài toán tìm kiếm toàn bộ cây đường ngắn nhất từ một điểm đến tất cả điểm khác (mặc dù Dijkstra có khả năng này).
\end{itemize}

\section{Phương pháp nghiên cứu}

\subsection{Phương pháp lý thuyết}

Sử dụng tài liệu học tập từ môn Nhập môn AI và các sách chuyên khảo (AIMA, CLRS) để nắm vững lý thuyết các thuật toán, bao gồm:
\begin{itemize}
  \item Đọc và phân tích pseudocode từ tài liệu.
  \item Hiểu rõ điều kiện admissibility của heuristic, tính hoàn chỉnh và tối ưu.
  \item Phân tích độ phức tạp thời gian và không gian.
\end{itemize}

\subsection{Phương pháp thực nghiệm}

Xây dựng ứng dụng và chạy các thí nghiệm:
\begin{itemize}
  \item Cài đặt đầy đủ 4 thuật toán bằng Python.
  \item Viết 29 test case bao gồm unit test (kiểm thử từng hàm) và integration test (kiểm thử toàn flow HTTP).
  \item Chạy benchmark trên những cặp ga khác nhau để đo explored node count, chi phí, thời gian thực thi.
  \item So sánh kết quả giữa các thuật toán.
\end{itemize}

\subsection{Phương pháp thiết kế}

Sử dụng quy trình thiết kế phần mềm chuẩn:
\begin{itemize}
  \item Phân tích yêu cầu (functional + non-functional).
  \item Thiết kế kiến trúc 3-tầng: Presentation, Business Logic, Data.
  \item Thiết kế API RESTful với JSON.
  \item Thiết kế giao diện người dùng (UI) responsive, sử dụng Leaflet cho bản đồ.
\end{itemize}

\section{Cấu trúc báo cáo}

Báo cáo gồm 5 chương chính:

\begin{description}
  \item[Chương 1 --- Mở đầu] Nêu lý do chọn đề tài, mục tiêu nghiên cứu, phạm vi, phương pháp, và cấu trúc báo cáo.

  \item[Chương 2 --- Cơ sở lý thuyết] Trình bày chi tiết bài toán tìm đường, biểu diễn đồ thị, phân tích 4 thuật toán (UCS, Greedy, A*, Dijkstra), hàm heuristic, và so sánh lý thuyết.

  \item[Chương 3 --- Phân tích và thiết kế] Liệt kê yêu cầu chức năng và phi chức năng, trình bày kiến trúc hệ thống, mô hình dữ liệu, thiết kế API, thiết kế giao diện.

  \item[Chương 4 --- Thực hiện và kết quả] Mô tả công nghệ sử dụng, cấu trúc mã nguồn, chi tiết cài đặt các thuật toán, kết quả kiểm thử, kết quả benchmark so sánh.

  \item[Chương 5 --- Kết luận] Tóm tắt những gì đạt được, nêu hạn chế hiện tại, và đề xuất hướng phát triển tương lai.
\end{description}

Ngoài ra, các phụ lục chứa hướng dẫn cài đặt và chạy ứng dụng, cũng như các tài liệu tham khảo.

% ===== CHƯƠNG 2: CƠ SỞ LÝ THUYẾT =====

\chapter{Cơ Sở Lý Thuyết}

\section{Bài toán tìm đường trên đồ thị có trọng số}

\subsection{Định nghĩa chính thức}

Cho một đồ thị vô hướng có trọng số $G = (V, E, w)$ với:
\begin{itemize}
  \item $V$ là tập hợp các đỉnh (vertices) --- trong bối cảnh này là các ga metro.
  \item $E$ là tập hợp các cạnh (edges) --- đoạn đường nối hai ga liền kề.
  \item $w: E \to \mathbb{R}^+$ là hàm trọng số gán cho mỗi cạnh một giá trị dương (chi phí).
\end{itemize}

Bài toán tìm đường ngắn nhất từ ga nguồn $s$ đến ga đích $g$ là tìm một đường đi $P = [s = v_0, v_1, \ldots, v_n = g]$ sao cho tổng chi phí:
\[
  \text{cost}(P) = \sum_{i=0}^{n-1} w(v_i, v_{i+1})
\]
là nhỏ nhất trong tất cả các đường đi từ $s$ đến $g$.

\subsection{Không gian trạng thái (State Space)}

Mỗi trạng thái trong quá trình tìm kiếm được định nghĩa là:
\[
  \text{State} = (\text{station\_name}, \text{cumulative\_cost}, \text{path}, \text{explored\_count})
\]

Trạng thái bắt đầu: $s_0 = (s, 0, [s], 0)$

Trạng thái mục tiêu: Bất kỳ trạng thái nào có station\_name = $g$

\subsection{Hàm hành động (Actions)}

Từ một trạng thái hiện tại, có thể chuyển tới các trạng thái lân cận bằng cách đi qua các cạnh từ ga hiện tại tới các ga hàng xóm.

\section{Biểu diễn đồ thị trong ứng dụng}

Ứng dụng sử dụng biểu diễn **danh sách kề** (adjacency list). Cấu trúc dữ liệu chính:

\begin{lstlisting}[language=Python, caption=Cấu trúc dữ liệu Station và Edge (metro\_app/data.py)]
@dataclass(frozen=True)
class Station:
    name: str
    lat: float
    lon: float

@dataclass(frozen=True)
class Edge:
    target: str
    travel_time: float
    line: str
    distance_km: float
\end{lstlisting}

Đồ thị được lưu dưới dạng:
\[
  \text{graph: dict[str, list[Edge]]}
\]

Với mỗi ga (key), lưu danh sách các cạnh tới các ga hàng xóm.

\section{Uniform Cost Search (UCS)}

\subsection{Ý tưởng chính}

UCS là một thuật toán tìm kiếm hướng mục tiêu (goal-directed) sử dụng hàng đợi ưu tiên (priority queue). Thuật toán luôn chọn trạng thái có chi phí nhỏ nhất để khám phá, đảm bảo rằng lần đầu tiên đạt được mục tiêu là lần với chi phí nhỏ nhất.

\subsection{Pseudocode}

\begin{lstlisting}[caption=Uniform Cost Search Pseudocode]
function UniformCostSearch(graph, start, goal):
  queue = PriorityQueue()
  queue.push((0, start))

  came_from = {start: None}
  best_cost = {start: 0}
  explored = 0

  while queue is not empty:
    current_cost, current = queue.pop()

    // Skip if already found better path
    if current_cost > best_cost[current]:
      continue

    explored += 1

    // Found goal
    if current == goal:
      path = reconstruct_path(came_from, goal)
      return (path, current_cost, explored)

    // Expand neighbors
    for edge in graph[current]:
      new_cost = current_cost + edge.travel_time
      if edge.target not in best_cost or new_cost < best_cost[edge.target]:
        best_cost[edge.target] = new_cost
        came_from[edge.target] = current
        queue.push((new_cost, edge.target))

  return FAILURE
\end{lstlisting}

\subsection{Tính chất}

\begin{itemize}
  \item \textbf{Hoàn chỉnh (Complete):} Có. Nếu tồn tại đường đi, UCS sẽ tìm thấy.
  \item \textbf{Tối ưu (Optimal):} Có. Lần đầu pop goal từ queue là lần với chi phí nhỏ nhất.
  \item \textbf{Heuristic:} Không dùng heuristic.
  \item \textbf{Độ phức tạp thời gian:} $O(b^{C^*/\epsilon})$ với $b$ là branching factor, $C^*$ là chi phí tối ưu, $\epsilon$ là chi phí cạnh nhỏ nhất.
  \item \textbf{Độ phức tạp không gian:} $O(b^{C^*/\epsilon})$ --- lưu tất cả node trong queue.
\end{itemize}

\section{Thuật toán Dijkstra}

\subsection{Ý tưởng chính}

Dijkstra là một trường hợp đặc biệt của UCS khi heuristic $h = 0$. Tuy nhiên, trong lý thuyết truyền thống, Dijkstra được triển khai khác biệt với mục tiêu tính toàn bộ **Single-Source Shortest Path (SSSP)** từ một nguồn đến tất cả các đỉnh khác, không dừng sớm khi tìm được một mục tiêu cụ thể.

\subsection{Pseudocode}

\begin{lstlisting}[caption=Dijkstra Algorithm Pseudocode]
function Dijkstra(graph, start, goal):
  distances = {node: infinity for all nodes}
  distances[start] = 0

  came_from = {start: None}
  unvisited = {all nodes}
  explored = 0

  while unvisited is not empty:
    current = node in unvisited with min distances[node]

    if distances[current] == infinity:
      break  // No path exists

    unvisited.remove(current)
    explored += 1

    // Continue to compute full SSSP (no early exit)
    for edge in graph[current]:
      neighbor = edge.target
      new_distance = distances[current] + edge.travel_time

      if new_distance < distances[neighbor]:
        distances[neighbor] = new_distance
        came_from[neighbor] = current

  // Reconstruct path to goal using came_from
  path = reconstruct_path(came_from, goal)
  return (path, distances[goal], explored)
\end{lstlisting}

\subsection{Tính chất}

\begin{itemize}
  \item \textbf{Hoàn chỉnh:} Có.
  \item \textbf{Tối ưu:} Có.
  \item \textbf{Heuristic:} Không.
  \item \textbf{Độ phức tạp thời gian:} $O((V + E) \log V)$ với priority queue.
  \item \textbf{Độ phức tạp không gian:} $O(V)$ --- lưu distances table.
\end{itemize}

\subsection{UCS vs Dijkstra}

UCS và Dijkstra đều tính toàn bộ chi phí tối ưu, nhưng có sự khác biệt:

\begin{itemize}
  \item \textbf{UCS:} Goal-directed. Dừng ngay khi pop mục tiêu từ queue. Số node explored tối thiểu.
  \item \textbf{Dijkstra:} SSSP-oriented. Tính toàn bộ distances từ một nguồn đến tất cả node. Số node explored = toàn bộ graph reachable từ start, thường > UCS.
\end{itemize}

Trong báo cáo này, Dijkstra được cài đặt với early-exit (dừng khi pop goal), nhưng đạt được SSSP trước khi reconstruct path, nên số node explored cao hơn UCS.

\section{Greedy Best-First Search}

\subsection{Ý tưởng chính}

Greedy chỉ dùng heuristic estimate (không dùng chi phí thực), luôn chọn node có $h(n)$ nhỏ nhất (gần mục tiêu nhất theo ước lượng). Thuật toán nhanh nhưng không đảm bảo tối ưu.

\subsection{Pseudocode}

\begin{lstlisting}[caption=Greedy Best-First Search Pseudocode]
function GreedyBestFirstSearch(graph, stations, start, goal):
  queue = PriorityQueue()
  h_start = heuristic(stations, start, goal)
  queue.push((h_start, 0, start))

  came_from = {start: None}
  visited = {set of visited nodes}
  explored = 0

  while queue is not empty:
    _, current_cost, current = queue.pop()

    if current in visited:
      continue

    visited.add(current)
    explored += 1

    if current == goal:
      path = reconstruct_path(came_from, goal)
      return (path, current_cost, explored)

    for edge in graph[current]:
      neighbor = edge.target
      if neighbor not in visited:
        h_neighbor = heuristic(stations, neighbor, goal)
        new_cost = current_cost + edge.travel_time
        came_from[neighbor] = current
        queue.push((h_neighbor, new_cost, neighbor))

  return FAILURE
\end{lstlisting}

\subsection{Tính chất}

\begin{itemize}
  \item \textbf{Hoàn chỉnh:} Có (với graph hữu hạn và tracking visited).
  \item \textbf{Tối ưu:} Không. Greedy có thể tìm thấy đường không phải chi phí nhỏ nhất.
  \item \textbf{Heuristic:} Có.
  \item \textbf{Độ phức tạp thời gian:} $O(b^m)$ với $m$ là độ sâu bài toán.
  \item \textbf{Độ phức tạp không gian:} $O(b^m)$.
\end{itemize}

\section{Thuật toán A*}

\subsection{Ý tưởng chính}

A* kết hợp chi phí thực $g(n)$ (từ start tới $n$) và ước lượng $h(n)$ (từ $n$ tới goal) qua công thức:
\[
  f(n) = g(n) + h(n)
\]

A* luôn khám phá node có $f(n)$ nhỏ nhất, đảm bảo tối ưu nếu $h(n)$ là admissible (không overestimate).

\subsection{Pseudocode}

\begin{lstlisting}[caption=A* Algorithm Pseudocode]
function AStar(graph, stations, start, goal):
  queue = PriorityQueue()
  g_start = 0
  h_start = heuristic(stations, start, goal)
  queue.push((g_start + h_start, g_start, start))

  came_from = {start: None}
  best_g = {start: 0}
  explored = 0

  while queue is not empty:
    _, g_current, current = queue.pop()

    if g_current > best_g[current]:
      continue

    explored += 1

    if current == goal:
      path = reconstruct_path(came_from, goal)
      return (path, g_current, explored)

    for edge in graph[current]:
      neighbor = edge.target
      g_neighbor = g_current + edge.travel_time

      if neighbor not in best_g or g_neighbor < best_g[neighbor]:
        best_g[neighbor] = g_neighbor
        came_from[neighbor] = current
        h_neighbor = heuristic(stations, neighbor, goal)
        f_neighbor = g_neighbor + h_neighbor
        queue.push((f_neighbor, g_neighbor, neighbor))

  return FAILURE
\end{lstlisting}

\subsection{Tính chất}

\begin{itemize}
  \item \textbf{Hoàn chỉnh:} Có.
  \item \textbf{Tối ưu:} Có (khi $h(n)$ admissible).
  \item \textbf{Heuristic:} Có.
  \item \textbf{Độ phức tạp thời gian:} $O(b^{C^*/\epsilon})$ --- có thể tốt hơn UCS nếu heuristic tốt.
  \item \textbf{Độ phức tạp không gian:} $O(b^{C^*/\epsilon})$.
\end{itemize}

\section{Hàm Heuristic và Tính Admissibility}

\subsection{Định nghĩa Admissibility}

Một heuristic $h(n)$ được gọi là **admissible** nếu nó không bao giờ overestimate chi phí thực từ $n$ tới goal:
\[
  h(n) \leq h^*(n) \quad \forall n
\]

với $h^*(n)$ là chi phí thực tế nhỏ nhất từ $n$ tới goal.

Nếu $h(n)$ admissible, A* được chứng minh là optimal.

\subsection{Hàm Heuristic trong ứng dụng}

Ứng dụng sử dụng:

\begin{lstlisting}[language=Python, caption=Hàm Heuristic (metro\_app/algorithms.py)]
def heuristic(stations: dict[str, Station], current: str, goal: str) -> float:
    here = stations[current]
    destination = stations[goal]
    return haversine_km(here.lat, here.lon, destination.lat, destination.lon) / 35.0 * 60.0
\end{lstlisting}

Công thức: $h(n) = \frac{\text{haversine}(n, goal)}{35} \times 60$

\subsection{Tính Admissibility}

Hàm này admissible vì:
\begin{enumerate}
  \item Haversine tính khoảng cách thẳng (đường tròn lớn) giữa hai điểm.
  \item Trong metro, đường đi thực tế luôn $\geq$ đường thẳng (do phải đi theo tuyến).
  \item Do đó: haversine $\leq$ khoảng cách thực tế.
  \item Chia cho 35 km/h (tốc độ metro) và nhân 60 (chuyển sang phút) cho ra giá trị time estimate luôn $\leq$ chi phí thực tế.
  \item Suy ra $h(n)$ admissible.
\end{enumerate}

\section{So sánh Lý Thuyết Bốn Thuật toán}

\begin{table}[H]
  \centering
  \caption{So sánh lý thuyết bốn thuật toán tìm đường}
  \label{tab:comparison}
  \begin{tabular}{lcccccc}
    \toprule
    \textbf{Tiêu chí} & \textbf{UCS} & \textbf{Greedy} & \textbf{A*} & \textbf{Dijkstra} \\
    \midrule
    Hoàn chỉnh (Complete) & ✓ & ✓* & ✓ & ✓ \\
    Tối ưu (Optimal) & ✓ & ✗ & ✓** & ✓ \\
    Dùng heuristic & ✗ & ✓ & ✓ & ✗ \\
    Time complexity & $O(b^{C^*/\epsilon})$ & $O(b^m)$ & $O(b^{C^*/\epsilon})$ & $O((V+E)\log V)$ \\
    Space complexity & $O(b^{C^*/\epsilon})$ & $O(b^m)$ & $O(b^{C^*/\epsilon})$ & $O(V)$ \\
    \bottomrule
  \end{tabular}
\end{table}

*Greedy hoàn chỉnh với graph hữu hạn và tracking visited.

**A* tối ưu khi $h(n)$ admissible.

\section{Kết luận chương}

Chương này đã trình bày chi tiết bốn thuật toán tìm đường, mỗi thuật toán có những ưu điểm và nhược điểm riêng. UCS và Dijkstra đảm bảo tối ưu nhưng không dùng heuristic. Greedy nhanh nhưng không tối ưu. A* cân bằng tốc độ và tính tối ưu nhờ vào heuristic admissible. Chương tiếp theo sẽ trình bày cách áp dụng những thuật toán này vào ứng dụng thực tế.

% ===== CHƯƠNG 3: PHÂN TÍCH VÀ THIẾT KẾ =====

\chapter{Phân Tích và Thiết Kế}

\section{Phân tích yêu cầu}

\subsection{Yêu cầu chức năng}

\begin{enumerate}
  \item \textbf{R1 --- Chọn ga từ danh sách:} Người dùng có thể chọn ga đi và ga đến từ dropdown danh sách 411 ga, có hỗ trợ typeahead/search nhanh.

  \item \textbf{R2 --- Chọn ga từ bản đồ:} Người dùng click nút "Điểm đi" hoặc "Điểm đến", sau đó click vị trí trên bản đồ để chọn điểm. Hệ thống tự tìm ga gần nhất (tối đa 50 km).

  \item \textbf{R3 --- Chọn thuật toán:} Người dùng chọn một trong 4 thuật toán (UCS, Greedy, A*, Dijkstra) từ dropdown.

  \item \textbf{R4 --- Cấm tuyến:} Người dùng có thể chọn cấm một hoặc nhiều tuyến metro trước khi chạy thuật toán.

  \item \textbf{R5 --- Cấm đoạn:} Người dùng có thể cấm một hoặc nhiều đoạn (cạnh) cụ thể giữa 2 ga.

  \item \textbf{R6 --- Tìm đường:} Hệ thống xử lý request tìm đường, gọi thuật toán tương ứng, trả lại kết quả.

  \item \textbf{R7 --- Hiển thị kết quả:} Hiển thị đường đi trên bản đồ bằng polyline, danh sách ga, quãng đường, chi phí, số node explored, danh sách tuyến.

  \item \textbf{R8 --- Bộ lọc hiển thị:} Popover cho phép toggle hiển thị/ẩn trạm và từng tuyến (visual only, không ảnh hưởng pathfinding).

  \item \textbf{R9 --- Nút "Tìm tuyến mới":} Reset route nhưng giữ nguyên ràng buộc đã chọn.

  \item \textbf{R10 --- Responsive design:} Giao diện hoạt động tốt trên desktop (1280+px), tablet, mobile (360px+).
\end{enumerate}

\subsection{Yêu cầu phi chức năng}

\begin{itemize}
  \item \textbf{NF1 --- Hiệu năng:} Response time API $< 50$ ms cho query điển hình.
  \item \textbf{NF2 --- Khả năng mở rộng:} Hỗ trợ tối thiểu 411 ga (current), tối đa 1000+ ga.
  \item \textbf{NF3 --- Độ tin cậy:} 99%+ uptime (với data hợp lệ).
  \item \textbf{NF4 --- Usability:} Giao diện trực quan, focus ring rõ ràng (accessibility AA).
  \item \textbf{NF5 --- Caching:} Polyline cache để tránh tính toán lại khi toggle blocked lines.
\end{itemize}

\section{Kiến trúc hệ thống}

Ứng dụng sử dụng kiến trúc **3-tầng** (Three-tier architecture):

% ===== PLACEHOLDER HÌNH =====
% TODO: Vẽ sơ đồ kiến trúc 3 tầng: Presentation Layer (Browser) ↔ Business Logic Layer (Service + Algorithms) ↔ Data Layer (Graph + JSON).
% Gợi ý file: figures/architecture-diagram.png
% Kích thước đề xuất: 0.8\textwidth
\begin{figure}[H]
  \centering
  \fbox{\begin{minipage}{0.8\textwidth}\centering\vspace{3cm}%
    \textit{[Placeholder: Kiến trúc 3-tầng hệ thống]}%
  \vspace{3cm}\end{minipage}}
  \caption{Kiến trúc hệ thống 3-tầng: Presentation, Business Logic, Data}
  \label{fig:architecture}
\end{figure}

\subsection{Tầng Presentation (Giao diện)}

- \textbf{Công nghệ:} HTML5, CSS3 (custom properties, Grid, Flexbox), JavaScript ES2020.
- \textbf{Component chính:}
  \begin{itemize}
    \item Dropdown chọn ga (station mode).
    \item Button chế độ pick điểm map.
    \item Dropdown chọn thuật toán.
    \item Checkbox cấm tuyến/đoạn.
    \item Leaflet map với polyline route.
    \item Popover filter (responsive).
    \item Panel hiển thị kết quả (path, distance, cost, explored).
  \end{itemize}
- \textbf{API Client:} Fetch API để gọi backend.

\subsection{Tầng Business Logic}

- \textbf{Công nghệ:} Python, HTTP server (stdlib).
- \textbf{Module chính:}
  \begin{itemize}
    \item \textbf{service.py}: Các hàm higher-level như `find_routes()`, `find_nearest_station()`, filter graph.
    \item \textbf{algorithms.py}: Cài đặt 4 thuật toán (UCS, Greedy, A*, Dijkstra).
    \item \textbf{geo.py}: Tính toán Haversine distance.
    \item \textbf{server.py}: HTTP routing, request/response parsing.
  \end{itemize}
- \textbf{Endpoints API:} `/api/network`, `/api/routes`, `/api/nearest-station`, `/api/routes-by-points`.

\subsection{Tầng Data}

- \textbf{Công nghệ:} JSON file, in-memory Python dict.
- \textbf{Dữ liệu chính:}
  \begin{itemize}
    \item \textbf{Station:} name, lat, lon (411 ga).
    \item \textbf{Edge:} source, target, line, travel\_time, distance\_km.
    \item \textbf{Graph:} Biểu diễn danh sách kề.
  \end{itemize}
- \textbf{Nguồn:} OpenStreetMap (file `data/shanghai_metro_osm.json`, 206 KB).

\section{Mô hình dữ liệu}

\begin{lstlisting}[language=Python, caption=Dataclass Station và Edge (metro\_app/data.py)]
from dataclasses import dataclass

@dataclass(frozen=True)
class Station:
    name: str
    lat: float
    lon: float

@dataclass(frozen=True)
class Edge:
    target: str
    travel_time: float
    line: str
    distance_km: float
\end{lstlisting}

Đồ thị được lưu: `graph: dict[str, list[Edge]]`

Ví dụ: `graph["Hongqiao Railway Station"] = [Edge("Songhong Road", 6.0, "Line 2", 4.8), ...]`

\section{Luồng xử lý request tìm đường}

% ===== PLACEHOLDER HÌNH =====
% TODO: Vẽ sequence diagram: Browser → Server → Service → Algorithms → return result → Browser display
% Gợi ý file: figures/sequence-diagram-find-route.png
% Kích thước đề xuất: 0.85\textwidth
\begin{figure}[H]
  \centering
  \fbox{\begin{minipage}{0.85\textwidth}\centering\vspace{3.5cm}%
    \textit{[Placeholder: Sequence diagram luồng xử lý tìm đường]}%
  \vspace{3.5cm}\end{minipage}}
  \caption{Sequence diagram: User click → API request → Algorithm → Response → Display result}
  \label{fig:sequence}
\end{figure}

Luồng chi tiết:
\begin{enumerate}
  \item Người dùng click button "Tìm đường" trên UI.
  \item JavaScript xây dựng request object (start, goal, algorithm, blocked\_lines, blocked\_segments).
  \item Gửi HTTP GET request tới `/api/routes-by-points` hoặc `/api/routes`.
  \item Server.py nhận request, parse parameters.
  \item Gọi `service.find_routes()` với các tham số.
  \item Service tạo graph đã lọc (loại bỏ tuyến/đoạn bị cấm).
  \item Service gọi hàm thuật toán tương ứng (UCS/Greedy/A*/Dijkstra).
  \item Thuật toán trả lại `SearchResult` (path, cost, explored\_count, line\_sequence, distance).
  \item Service format JSON response.
  \item Server trả lại HTTP 200 + JSON body.
  \item JavaScript parse response, vẽ polyline, hiển thị kết quả.
\end{enumerate}

\section{Thiết kế API}

\begin{table}[H]
  \centering
  \caption{Các API endpoint chính}
  \label{tab:api}
  \begin{tabular}{lll}
    \toprule
    \textbf{Endpoint} & \textbf{Phương thức} & \textbf{Mô tả} \\
    \midrule
    \texttt{/api/network} & GET & Metadata mạng (stations, edges, lines, algorithms) \\
    \texttt{/api/routes} & GET & Tìm đường bằng station names \\
    \texttt{/api/nearest-station} & GET & Tìm ga gần nhất từ lat/lon \\
    \texttt{/api/routes-by-points} & GET & Tìm đường từ lat/lon (tự map tới ga gần nhất) \\
    \bottomrule
  \end{tabular}
\end{table}

\section{Thiết kế giao diện người dùng}

UI có 3 chế độ chính:

\subsection{Chế độ "Chọn ga"}

Người dùng:
\begin{enumerate}
  \item Chọn "Ga đi" từ dropdown (411 ga, sorted, hỗ trợ typeahead).
  \item Chọn "Ga đến" từ dropdown.
  \item (Tuỳ chọn) Tick "Tuyến bị cấm" để loại tuyến.
  \item (Tuỳ chọn) Chọn 2 ga trong phần "Đoạn bị cấm", bấm "Thêm đoạn".
  \item Chọn "Thuật toán".
  \item Bấm "Tìm đường".
\end{enumerate}

% ===== PLACEHOLDER HÌNH =====
% TODO: Screenshot chế độ chọn ga, hiển thị dropdown ga đi, ga đến, dropdown thuật toán, checkbox cấm tuyến
% Gợi ý file: figures/ui-station-mode.png
% Kích thước đề xuất: 0.75\textwidth
\begin{figure}[H]
  \centering
  \fbox{\begin{minipage}{0.75\textwidth}\centering\vspace{3cm}%
    \textit{[Placeholder: Screenshot chế độ chọn ga]}%
  \vspace{3cm}\end{minipage}}
  \caption{Giao diện chế độ "Chọn ga" với dropdown, checkbox cấm tuyến}
  \label{fig:ui-station}
\end{figure}

\subsection{Chế độ "Click trên map"}

Người dùng:
\begin{enumerate}
  \item Bấm nút "Điểm đi" → cursor → click map.
  \item Hệ thống tìm ga gần nhất, lưu tọa độ.
  \item Tương tự với "Điểm đến".
  \item Bấm "Tìm đường từ điểm đã chọn".
\end{enumerate}

% ===== PLACEHOLDER HÌNH =====
% TODO: Screenshot chế độ map mode với popover filter bên phải, thể hiện toàn bộ metro trên bản đồ, điểm pick được hiển thị
% Gợi ý file: figures/ui-map-mode.png
% Kích thước đề xuất: 0.8\textwidth
\begin{figure}[H]
  \centering
  \fbox{\begin{minipage}{0.8\textwidth}\centering\vspace{3.5cm}%
    \textit{[Placeholder: Screenshot chế độ map mode với filter popover]}%
  \vspace{3.5cm}\end{minipage}}
  \caption{Giao diện map mode: Leaflet map, filter popover (checkbox tuyến), highlight route}
  \label{fig:ui-map}
\end{figure}

\subsection{Popover Filter}

Bấm nút "⚙ Hiển thị" để mở popover:
\begin{itemize}
  \item Checkbox "Trạm" --- toggle marker ga.
  \item Checkbox "Tất cả tuyến" (master) --- toggle hàng loạt.
  \item 20 checkbox từng tuyến (Line 1, Line 2, ..., Line 17) --- toggle riêng.
\end{itemize}

\textbf{Lưu ý:} Filter này chỉ ảnh hưởng visual, không đổi kết quả pathfinding.

\section{Kết luận chương}

Chương này đã mô tả chi tiết thiết kế hệ thống, từ yêu cầu, kiến trúc 3-tầng, mô hình dữ liệu, luồng xử lý, API, đến giao diện người dùng. Những thiết kế này đảm bảo ứng dụng có thể mở rộng, dễ bảo trì, và cung cấp trải nghiệm người dùng tốt trên nhiều thiết bị.

% ===== CHƯƠNG 4: THỰC HIỆN VÀ KẾT QUẢ =====

\chapter{Thực Hiện và Kết Quả}

\section{Công nghệ sử dụng}

\subsection{Backend}

\begin{itemize}
  \item \textbf{Ngôn ngữ:} Python 3.10+ (sử dụng type hints, dataclass, walrus operator).
  \item \textbf{Thư viện:} Chỉ dùng Python Standard Library:
    \begin{itemize}
      \item \texttt{http.server} --- HTTP server threading.
      \item \texttt{urllib.parse} --- URL parsing.
      \item \texttt{json} --- JSON serialization.
      \item \texttt{math} --- Tính toán Haversine.
      \item \texttt{heapq} --- Priority queue cho thuật toán.
      \item \texttt{dataclasses} --- Station, Edge.
      \item \texttt{pathlib} --- Path handling.
    \end{itemize}
  \item \textbf{Database:} Không dùng. Dữ liệu lưu trong JSON file và load vào memory.
\end{itemize}

\subsection{Frontend}

\begin{itemize}
  \item \textbf{HTML5:} Semantic markup.
  \item \textbf{CSS3:} Custom properties, CSS Grid, Flexbox, media queries (responsive).
  \item \textbf{JavaScript:} ES2020 (const/let, arrow functions, template literals, spread operator, async/await).
  \item \textbf{Leaflet.js v1.9.4:} Interactive web map, markers, polylines, Popup.
  \item \textbf{OpenStreetMap:} Tile layer (không cần API key).
\end{itemize}

\subsection{Deployment \& Testing}

\begin{itemize}
  \item \textbf{Testing Framework:} Python \texttt{unittest} (built-in).
  \item \textbf{Deployment:} Standalone Python HTTP server (không cần app server như Gunicorn, nếu volume tăng có thể thêm).
  \item \textbf{Containerization:} (Optional) Docker.
\end{itemize}

\section{Cấu trúc mã nguồn}

\begin{lstlisting}[caption=Cấu trúc thư mục project]
Project-for-Introduction-to-AI/
|-- app.py                          # Entry point
|-- server.py                       # HTTP server + endpoints
|-- metro_app/
|   |-- __init__.py
|   |-- data.py                     # Station, Edge, load_network()
|   |-- algorithms.py               # UCS, Greedy, A*, Dijkstra
|   |-- service.py                  # Business logic (find_routes, etc.)
|   |-- geo.py                      # haversine_km()
|   `-- osm_import.py               # (Optional) Refresh data from Overpass API
|-- web/
|   |-- index.html
|   |-- app.js
|   `-- styles.css
|-- data/
|   `-- shanghai_metro_osm.json     # 411 ga, 20 tuyen (206 KB)
|-- tests/
|   |-- test_algorithms.py
|   `-- test_http_server.py
`-- README.md
\end{lstlisting}

\section{Thực hiện 4 thuật toán}

\subsection{Uniform Cost Search}

\begin{lstlisting}[language=Python, caption=UCS implementation (metro\_app/algorithms.py - excerpt)]
def uniform_cost_search(
    graph: dict[str, list[Edge]],
    start: str,
    goal: str,
) -> SearchResult:
    queue: list[tuple[float, int, str]] = [(0.0, 0, start)]
    tracker = count(1)
    came_from: dict[str, tuple[str, str] | None] = {start: None}
    best_cost = {start: 0.0}
    explored = 0

    while queue:
        current_cost, _, current = heapq.heappop(queue)
        if current_cost > best_cost[current]:
            continue

        explored += 1
        if current == goal:
            path, lines = _reconstruct_path(came_from, goal)
            return SearchResult("UCS", path, round(current_cost, 2), explored,
                                lines, _path_distance(graph, path, lines))

        for edge in graph[current]:
            new_cost = current_cost + edge.travel_time
            if edge.target not in best_cost or new_cost < best_cost[edge.target]:
                best_cost[edge.target] = new_cost
                came_from[edge.target] = (current, edge.line)
                heapq.heappush(queue, (new_cost, next(tracker), edge.target))

    raise ValueError(f"Cannot reach {goal} from {start}.")
\end{lstlisting}

\subsection{Dijkstra}

Dijkstra được cài đặt tương tự UCS nhưng tính toàn bộ SSSP trước khi reconstruct path, dẫn tới số node explored lớn hơn.

\subsection{Greedy Best-First Search}

\begin{lstlisting}[language=Python, caption=Greedy BFS pseudocode (thực hiện tương tự)]
def greedy_best_first_search(
    graph: dict[str, list[Edge]],
    stations: dict[str, Station],
    start: str,
    goal: str,
) -> SearchResult:
    # Dùng heuristic h(n) làm priority
    # Early-exit khi pop goal từ queue
    pass
\end{lstlisting}

\subsection{A* Search}

\begin{lstlisting}[language=Python, caption=A* pseudocode]
def a_star_search(
    graph: dict[str, list[Edge]],
    stations: dict[str, Station],
    start: str,
    goal: str,
) -> SearchResult:
    # Dùng f(n) = g(n) + h(n) làm priority
    # g(n) = chi phí từ start, h(n) = heuristic
    pass
\end{lstlisting}

\section{Hàm Heuristic}

\begin{lstlisting}[language=Python, caption=Hàm heuristic (metro\_app/algorithms.py)]
def heuristic(stations: dict[str, Station], current: str, goal: str) -> float:
    here = stations[current]
    destination = stations[goal]
    return haversine_km(here.lat, here.lon, destination.lat, destination.lon) / 35.0 * 60.0
\end{lstlisting}

Công thức: Khoảng cách Haversine (km) ÷ 35 (km/h) × 60 (phút).

\begin{lstlisting}[language=Python, caption=Hàm Haversine (metro\_app/geo.py)]
def haversine_km(lat1: float, lon1: float, lat2: float, lon2: float) -> float:
    from math import radians, sin, cos, sqrt, atan2

    R = 6371  # Earth radius in km

    phi1, phi2 = radians(lat1), radians(lat2)
    delta_phi = radians(lat2 - lat1)
    delta_lambda = radians(lon2 - lon1)

    a = sin(delta_phi/2)**2 + cos(phi1) * cos(phi2) * sin(delta_lambda/2)**2
    c = 2 * atan2(sqrt(a), sqrt(1-a))

    return R * c
\end{lstlisting}

\section{Kiểm thử}

Bộ test gồm 29 test case:

\subsection{Unit Tests (metro\_app tests)}

\begin{itemize}
  \item \textbf{14 tests thuật toán:} Mỗi thuật toán được test trên các đồ thị sample khác nhau (simple path, multiple paths, branching).
  \item \textbf{3 tests correctness:} Verify kết quả path, cost, explored count so sánh với expected.
  \item \textbf{2 tests Dijkstra:} So sánh SSSP đầy đủ vs goal-directed behavior.
  \item \textbf{Test filter:} Verify graph filter khi loại tuyến/đoạn.
  \item \textbf{Test geo bounds:} Verify nearest-station reject điểm > 50 km.
\end{itemize}

\subsection{Integration Tests (HTTP server tests)}

\begin{itemize}
  \item \textbf{4 endpoints × (happy path + malformed input):} Test \texttt{/api/network}, \texttt{/api/routes}, \texttt{/api/nearest-station}, \texttt{/api/routes-by-points}.
  \item \textbf{10 tests HTTP:} Gồm status code 200, 400, content-type JSON.
\end{itemize}

\begin{lstlisting}[language=bash, caption=Chạy kiểm thử]
$ python -m unittest discover -s tests -v
test_algorithms.py ............. [19 tests]
test_http_server.py ............ [10 tests]
Ran 29 tests in 0.6s
OK
\end{lstlisting}

\section{Kết quả Benchmark}

Chạy benchmark trên một cặp điểm: **Hongqiao Railway Station → Century Avenue**

% ===== PLACEHOLDER BẢNG =====
% TODO: Điền dữ liệu benchmark thực tế từ project
% Bảng so sánh cost, distance, explored nodes, path length
\begin{table}[H]
  \centering
  \caption{Kết quả benchmark so sánh 4 thuật toán}
  \label{tab:benchmark}
  \begin{tabular}{lrrrr}
    \toprule
    \textbf{Thuật toán} & \textbf{Cost (phút)} & \textbf{Distance (km)} & \textbf{Explored} & \textbf{Path Length} \\
    \midrule
    UCS & 48.61 & 21.82 & 170 & 9 ga \\
    Greedy & 48.85 & 21.96 & 15 & 9 ga \\
    A* & 48.61 & 21.82 & 64 & 9 ga \\
    Dijkstra & 48.61 & 21.82 & $\geq 170$ & 9 ga \\
    \bottomrule
  \end{tabular}
\end{table}

\subsection{Phân tích kết quả}

\begin{itemize}
  \item \textbf{UCS, A*, Dijkstra:} Tất cả đều tối ưu, chi phí bằng nhau (48.61 phút, 21.82 km).

  \item \textbf{Greedy:} Chi phí cao hơn (48.85 phút, 21.96 km) vì chỉ dùng heuristic mà không dùng chi phí thực tế.

  \item \textbf{Explored nodes:}
    \begin{itemize}
      \item Greedy: 15 (nhanh nhất, nhưng không tối ưu).
      \item A*: 64 (cân bằng tốt, tối ưu).
      \item UCS: 170 (không dùng heuristic, nhưng tối ưu).
      \item Dijkstra: $\geq 170$ (tính toàn bộ SSSP).
    \end{itemize}

  \item \textbf{Path length:} Bốn thuật toán đều tìm path 9 ga (do tính tối ưu hoặc sự trùng lặp đường đi).
\end{itemize}

% ===== PLACEHOLDER HÌNH =====
% TODO: Vẽ biểu đồ bar chart so sánh explored\_nodes của 4 algo
% Gợi ý file: figures/benchmark-chart-explored-nodes.png
% Kích thước đề xuất: 0.75\textwidth
\begin{figure}[H]
  \centering
  \fbox{\begin{minipage}{0.75\textwidth}\centering\vspace{2.5cm}%
    \textit{[Placeholder: Biểu đồ bar chart so sánh explored nodes]}%
  \vspace{2.5cm}\end{minipage}}
  \caption{Biểu đồ so sánh số node explored bởi 4 thuật toán (Greedy=15, A*=64, UCS=170, Dijkstra=$\geq 170$)}
  \label{fig:benchmark-chart}
\end{figure}

\section{Đánh giá UX/UI}

\subsection{Chế độ chọn ga}

Ưu điểm:
\begin{itemize}
  \item Dropdown danh sách 411 ga, sorted alphabet, dễ tìm.
  \item Hỗ trợ typeahead/search (gõ chữ cái đầu nhảy nhanh).
  \item Checkbox cấm tuyến/đoạn trực quan, không cần đọc hướng dẫn.
\end{itemize}

Hạn chế:
\begin{itemize}
  \item Dropdown rất dài (411 item), người dùng phải scroll, có thể khó sử dụng trên mobile.
  \item Một số ga tên chỉ là số (ví dụ \textquotedbl 1\textquotedbl, \textquotedbl 2\textquotedbl) do data OSM thiếu tag \texttt{name}.
\end{itemize}

\subsection{Chế độ click map}

Ưu điểm:
\begin{itemize}
  \item Trực quan, người dùng không cần nhớ tên ga.
  \item Cursor đổi thành crosshair, rõ ràng là chế độ pick.
  \item Tự động thoát mode sau khi pick.
\end{itemize}

Hạn chế:
\begin{itemize}
  \item Tìm ga gần nhất dùng Haversine thẳng, không tính đường bộ thực tế.
  \item Ngưỡng 50 km có thể quá lớn/nhỏ tuỳ khu vực.
\end{itemize}

\subsection{Bộ lọc hiển thị}

\begin{itemize}
  \item Toggle popover bằng nút "⚙ Hiển thị" ở góc trên-phải map.
  \item 20 checkbox tuyến, hỗ trợ master checkbox "Tất cả tuyến".
  \item Visual only, không ảnh hưởng pathfinding (để tránh confusion).
\end{itemize}

\subsection{Responsive Design}

\begin{itemize}
  \item Desktop (1280+): Sidebar form bên trái, map toàn màn hình bên phải.
  \item Tablet (768-1279): Stacked layout, form trên, map dưới.
  \item Mobile (360-767): Fullscreen form/map, collapse section.
  \item Touch-friendly: Button min-height 40px, 44px primary button.
\end{itemize}

\subsection{Accessibility}

\begin{itemize}
  \item ARIA labels trên button (aria-pressed, aria-expanded, aria-controls).
  \item Focus ring rõ ràng (focus-visible, custom --shadow-focus).
  \item Contrast AA trên body text (#14181c trên #ffffff ~ 17:1).
  \item prefers-reduced-motion support.
  \item Semantic HTML (button, form, role="dialog" trên popover).
\end{itemize}

\section{Performance Metrics}

\begin{itemize}
  \item \textbf{API response:} $< 20$ ms cho query điển hình (Hongqiao → Century Avenue).
  \item \textbf{Test suite:} 29 tests, $\approx 0.6$ s.
  \item \textbf{Initial load:} HTML + JS + CSS $< 100$ KB (gzip).
  \item \textbf{Data load:} JSON (206 KB) parse vào memory $< 50$ ms.
  \item \textbf{Polyline rendering:} Cache enabled, toggle blocked-line setStyle() không rebuild $\approx 10$ × tốc độ.
\end{itemize}

\section{Kết luận chương}

Chương này đã mô tả chi tiết cài đặt ứng dụng, từ công nghệ sử dụng, cấu trúc code, thực hiện 4 thuật toán, kết quả kiểm thử, benchmark, đến đánh giá UX/UI. Kết quả benchmark cho thấy A* là sự cân bằng tốt giữa tốc độ (64 node explored) và tính tối ưu. Greedy nhanh nhất (15 node) nhưng không tối ưu. UCS/Dijkstra đảm bảo tối ưu nhưng không dùng heuristic nên explored nhiều node hơn A*.

% ===== CHƯƠNG 5: KẾT LUẬN =====

\chapter{Kết Luận}

\section{Những gì đã đạt được}

Báo cáo này đã hoàn thành các mục tiêu chính sau:

\subsection{Khía cạnh lý thuyết}

\begin{enumerate}
  \item \textbf{Nắm vững 4 thuật toán:} Hiểu rõ nguyên lý hoạt động, pseudocode, độ phức tạp, tính hoàn chỉnh, tính tối ưu của Uniform Cost Search, Greedy Best-First Search, A*, và Dijkstra.

  \item \textbf{Phân tích heuristic admissibility:} Chứng minh rằng hàm heuristic dựa trên Haversine / 35 km/h là admissible, do đó A* được đảm bảo tối ưu.

  \item \textbf{So sánh lý thuyết:} Xây dựng bảng so sánh 4 thuật toán theo nhiều tiêu chí (complete, optimal, time/space complexity, heuristic usage).
\end{enumerate}

\subsection{Khía cạnh ứng dụng thực tế}

\begin{enumerate}
  \item \textbf{Xây dựng ứng dụng web hoàn chỉnh:} Cài đặt 4 thuật toán bằng Python, xây dựng HTTP API RESTful với 4 endpoints, frontend responsive với Leaflet.js.

  \item \textbf{Làm việc với dữ liệu thực:} Sử dụng 411 ga và 20 tuyến metro Thượng Hải từ OpenStreetMap, xử lý dữ liệu địa lý (Haversine), estimate chi phí edge từ khoảng cách và tốc độ.

  \item \textbf{Kiến trúc 3-tầng:} Thiết kế và triển khai kiến trúc Presentation-Business Logic-Data, đảm bảo code modular, dễ bảo trì.

  \item \textbf{Giao diện người dùng:} Hai chế độ chọn ga (dropdown + click map), bộ lọc popover, responsive design (desktop/tablet/mobile), accessibility AA.
\end{enumerate}

\subsection{Khía cạnh kỹ thuật}

\begin{enumerate}
  \item \textbf{Viết test toàn diện:} 29 test case (14 + 3 + 2 + 10), bao gồm unit test thuật toán, correctness test, integration test HTTP.

  \item \textbf{Benchmark so sánh:} Chạy 4 thuật toán trên cùng cặp điểm, đo explored node count, chi phí, distance, rút ra những phát hiện về hiệu năng.

  \item \textbf{Publish codebase:} Code được organize rõ ràng, naming convention consistent, docstring rõ ràng, sẵn sàng cho collaboration.

  \item \textbf{Performance optimization:} Cache polyline để tránh tính toán lại khi toggle blocked lines ($\approx 10 \times$ speedup).
\end{enumerate}

\section{Hạn chế hiện tại}

\subsection{Hạn chế về dữ liệu}

\begin{enumerate}
  \item \textbf{Chi phí ước lượng, không chính xác 100\%:} OSM không cung cấp thời gian chạy thực tế từng đoạn, nên chi phí được ước lượng từ công thức (khoảng cách / 35 km/h + phụ phí). Kết quả "tối ưu" là tối ưu **với mô hình này**, không nhất thiết tối ưu trong thực tế.

  \item \textbf{Data quality issue:} Một số ga tên chỉ là số (ví dụ "1", "2") do OSM relation thiếu tag `name`. Một số ga có typo (ví dụ "National Exhibiation" thay vì "Exhibition").

  \item \textbf{Không tính giờ cao điểm, đóng cửa tuyến, transfer time:} Mô hình giả định tốc độ metro đều là 35 km/h và có sẵn chuyển tuyến, không tính thực tế.
\end{enumerate}

\subsection{Hạn chế về thuật toán}

\begin{enumerate}
  \item \textbf{Dijkstra tính SSSP:} Triển khai Dijkstra tính toàn bộ distances từ start tới tất cả node, không dừng sớm khi tìm goal, dẫn tới explored count cao hơn UCS.

  \item \textbf{Heuristic chỉ dựa trên khoảng cách:} Haversine là heuristic tổng quát nhưng không phải tối ưu nhất cho metro (không tính cấu trúc tuyến).
\end{enumerate}

\subsection{Hạn chế về giao diện}

\begin{enumerate}
  \item \textbf{Dropdown 411 ga rất dài:} Trên mobile, dropdown đôi khi khó sử dụng. Có thể cải thiện bằng search modal riêng.

  \item \textbf{Nearest station dùng Haversine thẳng:} Không tính đường bộ, đường sông, hay địa hình thực tế từ điểm click tới ga.

  \item \textbf{Không persist state:} F5 mất hết lựa chọn (blocked lines, picked points). Có thể lưu vào localStorage.

  \item \textbf{Không có walking directions:} Không hiển thị hướng đi chi tiết (đi thẳng tới ga, xuống cầu thang, v.v.).
\end{enumerate}

\section{Hướng phát triển tương lai}

\subsection{Cải thiện dữ liệu}

\begin{enumerate}
  \item \textbf{Lấy chi phí thực tế:} Nếu có API cung cấp thời gian chạy thực tế (ví dụ GTFS real-time data), có thể replace công thức hiện tại.

  \item \textbf{Thêm transfer time:} Ước lượng thời gian chuyển tuyến dựa trên khoảng cách giữa hai ga gần nhau trên các tuyến khác nhau.

  \item \textbf{Dữ liệu giờ cao điểm:} Nhập dữ liệu tốc độ khác nhau tuỳ theo giờ trong ngày.
\end{enumerate}

\subsection{Cải thiện thuật toán}

\begin{enumerate}
  \item \textbf{Cài đặt lại Dijkstra với early-exit:} Tính SSSP nhưng dừng ngay khi pop goal (hiện tại tính đầy đủ rồi mới reconstruct).

  \item \textbf{Heuristic pattern database:} Precompute heuristic values từ một subset ga, tăng accuracy.

  \item \textbf{Bidirectional search:} Tìm kiếm từ cả start và goal, gặp giữa đường, giảm không gian tìm kiếm.

  \item \textbf{Thêm thuật toán khác:} Bellman-Ford (xử lý edge weight âm nếu có), Floyd-Warshall (all-pairs shortest path), Contraction Hierarchies (tối ưu để query nhanh).
\end{enumerate}

\subsection{Cải thiện giao diện}

\begin{enumerate}
  \item \textbf{Search modal thay dropdown:} Cải thiện UX trên mobile, hỗ trợ fuzzy search station name.

  \item \textbf{Persist state:} Lưu lựa chọn vào localStorage hoặc URL query string, F5 vẫn giữ nguyên.

  \item \textbf{Walking directions:} Tích hợp OSRM hoặc Mapbox Directions API để cung cấp hướng đi chi tiết từ điểm user click tới ga.

  \item \textbf{Compare multiple routes:} Hiển thị cùng lúc kết quả của 4 thuật toán, giúp user so sánh.

  \item \textbf{Share URL:} Tạo URL ngắn gọn với tham số start, goal, algorithm, dễ chia sẻ.

  \item \textbf{Dark mode:} Hỗ trợ prefers-color-scheme.
\end{enumerate}

\subsection{Cải thiện kiểm thử và deployment}

\begin{enumerate}
  \item \textbf{End-to-end test:} Selenium/Playwright test giao diện browser.

  \item \textbf{Load test:} Apache JMeter hoặc Locust, test performance với 100+ concurrent user.

  \item \textbf{Docker + CI/CD:} GitHub Actions tự động run test, build Docker image, deploy lên cloud (AWS/GCP/Heroku).

  \item \textbf{Monitoring \& logging:} Thêm ELK stack (Elasticsearch, Logstash, Kibana) để monitor production.
\end{enumerate}

\subsection{Cải thiện tài liệu}

\begin{enumerate}
  \item \textbf{API documentation:} OpenAPI/Swagger specification.

  \item \textbf{Architecture decision records (ADR):} Ghi lại lý do chọn mỗi công nghệ, cấu trúc.

  \item \textbf{Contributing guide:} Hướng dẫn cho contributor muốn fork + PR.

  \item \textbf{Video demo:} Screencast hướng dẫn cách dùng ứng dụng.
\end{enumerate}

\section{Cảm ơn}

Tác giả xin cảm ơn \gvname{} đã hướng dẫn chi tiết trong suốt quá trình thực hiện bài tập lớn này. Cảm ơn các giáo viên khác trong Khoa Công Nghệ Thông Tin đã cung cấp tài liệu tham khảo quý báu.

Cảm ơn cộng đồng nguồn mở --- OpenStreetMap (dữ liệu), Leaflet.js (bản đồ), Python Standard Library (backend) --- đã làm cho dự án này có thể hoàn thành.

Cảm ơn những bạn cùng lớp đã trao đổi ý kiến, giúp đỡ em trong quá trình học tập.

\section{Tài liệu tham khảo}

Xem phần "Tài liệu tham khảo" ở cuối báo cáo.


% ===== BIBLIOGRAPHY =====
\cleardoublepage
\printbibliography[title=Tài liệu tham khảo]

% ===== APPENDIX =====
\appendix

\chapter{Hướng dẫn Cài đặt và Chạy}

\section{Yêu cầu Hệ thống}

\begin{itemize}
  \item Python 3.10+
  \item Trình duyệt hiện đại (Chrome, Firefox, Safari, Edge)
  \item Kết nối Internet (cho OpenStreetMap tiles)
  \item Không cần cài đặt thêm package --- chỉ dùng Python stdlib
\end{itemize}

\section{Hướng dẫn Chạy}

Để khởi động ứng dụng:

\begin{lstlisting}[language=bash, caption=Khởi động ứng dụng]
git clone https://github.com/3T216/Project-for-Introduction-to-AI
cd Project-for-Introduction-to-AI
python app.py
\end{lstlisting}

Sau đó mở trình duyệt tại: \url{http://127.0.0.1:8000}

\section{Chạy Kiểm Thử}

\begin{lstlisting}[language=bash, caption=Chạy bộ test]
python -m unittest discover -s tests -v
\end{lstlisting}

Điều này sẽ chạy 29 test case (14 tests thuật toán + 3 correctness + 2 Dijkstra + 10 HTTP).

\end{document}
